\section{What is tested, and by which rule}\label{sec:question}

\paragraph{The analysis.} Every transcript of the archive is cut into passages of 500 words, and a language
model answers a fixed bank of questions about the wording of each passage. A rating question has five described
levels, 0 to 4; the model's score is the mean of its own distribution over the levels. One model, \jev{}, rated
every passage; a second model from another maker, \drex{}, rated a large part again. Related questions are
grouped into eight dimensions (Table~\ref{tab:s2markers}), and a recording's score on a dimension is computed
from its passages. The site's methods page describes this in full. Agreement between two models shows that a
question can be applied consistently. It does not show that either model reads a passage as a person would.
That is what these studies test.

\paragraph{The comparison.} People read an excerpt and answer the same questions with the same five levels,
shown exactly as the models receive them, with ``cannot judge'' as a sixth answer. Both models rated exactly the
text people read: the first 250 words of a passage, as plain text without speaker labels. For one question, let
$y_{ir}$ be the rating of excerpt $i$ by rater $r$,
\begin{equation}\label{eq:measurement}
  y_{ir} = T_i + \varepsilon_{ir},
\end{equation}
where $T_i$ is what the raters of the population would say about the excerpt on average, the \emph{panel
consensus}, and $\varepsilon_{ir}$ everything that makes one rater differ from it. Two quantities carry the
verdicts:
\begin{align}
  \R{k} &= \frac{\sT}{\sT + \se/k}, \qquad \sT = \var(T),\ \se = \var(\varepsilon), && \text{the reliability of the mean of $k$ raters;} \label{eq:rel}\\
  r &= \rho\,\sqrt{\R{5}}, \qquad \rho = \corr(\text{score}, T), && \text{a model's correlation with the mean of five.} \label{eq:r}
\end{align}
$\R{5}$ is also the correlation to expect between the means of two separate panels of five. The second equation
separates what a correlation with a mean of five mixes: how well the model follows people ($\rho$) and how much
five people agree ($\R{5}$).

\paragraph{The rule.} Fixed before any rating was collected, with $r^\ast$ the correlation of the better of the
two models:
\begin{center}\small
\begin{tabular}{@{}ll@{}}
\toprule
Verdict & Rule \\
\midrule
\verd{usable} & $\R{5} \ge 0.60$ and $r^\ast \ge 0.50$ \\
\verd{unclear rubric} & $\R{5} < 0.40$: raters do not agree, so the wording of the question is the problem \\
\verd{misread} & $\R{5} \ge 0.60$ and $r^\ast < 0.30$: raters agree and the models do not follow them \\
\verd{borderline} & anything else \\
\bottomrule
\end{tabular}
\end{center}
The lines are conventions, not natural boundaries. A verdict is called \verd{settled} when the whole 95\%
interval of both quantities lies on one side of every line the rule compares it with; otherwise it is reported
as \verd{at the line}, and the estimate with its interval is the result.

\paragraph{The population.} A random 500-word passage of a random recording dated 2015 to 2025, videos and
livestreams in equal numbers, cut to its first 250 words. A recording of 40 minutes or more, or one filed as a
livestream, counts as a stream.

\paragraph{Fixed in advance, and added afterwards.} For each study the questions, the excerpts, the sets, the
checks that decide whose ratings count, the estimators and the rule were written down and tested on simulated
studies before collection. Those results are in Sections~\ref{sec:firstresults} and~\ref{sec:secondresults}. What
was computed after the results had been read is in Sections~\ref{sec:firstadded} and~\ref{sec:further} and is
marked as such.
